% Generated draft: requires mathematical and source review.
\documentclass[11pt]{article}
\usepackage[margin=25mm]{geometry}
\usepackage{fontspec}
\setmainfont{Latin Modern Roman}
\usepackage{amsmath,amssymb,mathtools,unicode-math}
\setmathfont{Latin Modern Math}
\usepackage{enumitem,needspace,xurl,hyperref}
\setlength{\parindent}{0pt}
\setlength{\parskip}{5pt}
\setlength{\emergencystretch}{3em}
\newcommand{\N}{\mathbb{N}}
\newcommand{\Z}{\mathbb{Z}}
\newcommand{\Q}{\mathbb{Q}}
\newcommand{\R}{\mathbb{R}}
\newcommand{\C}{\mathbb{C}}
\newcommand{\F}{\mathbb{F}}
\newcommand{\A}{\mathbb{A}}
\newcommand{\PP}{\mathbb{P}}
\newcommand{\E}{\mathbb{E}}
\newcommand{\eps}{\varepsilon}
\newcommand{\dd}{\,\mathrm{d}}
\begin{document}
\section*{UnsolvedMath ID: 30004695}
\textbf{Source problem code:} catalogue entry
\subsection{Definitions and mathematical statement}
The supplied record poses a comparative historiographical inquiry about mathematical works regarded as classics. It supplies no formal definition of a \emph{mathematical classic}, no specified collection of works or scholarly traditions, and no chronological or geographical boundaries for the expression \emph{Arabic world}. These expressions therefore retain their unresolved historical meanings; they are not defined here as mathematical predicates or domains.

The source question is:
\begin{quote}
What do mathematical classics mean in different scholarly traditions, how were they made available and used differently, and what did it mean for works to be deemed classics in the Arabic world?
\end{quote}
Its three related targets are to investigate the meanings of mathematical classics in different scholarly traditions, differences in their availability and use, and the significance of regarding works as classics in the Arabic world. The record does not specify the traditions, works, periods, evidence, or criteria for comparison. Consequently, it does not determine a quantified mathematical proposition or a unique formal statement. A more precise formulation requires historical scope and interpretive criteria that cannot be supplied without adding unsupported assumptions.

The registry labels this inquiry \emph{open}. Its supplied assessment describes active historiographical research questions rather than solved propositions; this is a source qualification, not an independently verified mathematical status.
\subsection{Short English statement}
Investigate how different scholarly traditions understood, made available, and used mathematical classics, including what classic status meant in the Arabic world. The source supplies a historical research question whose scope and key terms remain unspecified, so it does not support a precise mathematical proposition.
\subsection{Sources}
\begin{itemize}
\item Local registry: \texttt{data/unsolvedmath/problems.json}, supplied record with permanent UnsolvedMath ID \(30004695\), title \emph{Mathematical Classics across Scholarly Traditions}, source problem code \texttt{catalogue entry}, in the Oberwolfach Reports Open Problems collection.
\item Supplied report citation: \emph{Mathematics and its Ancient Classics Worldwide: Translations, Appropriations, Reconstructions, Roles} (2022), \emph{Oberwolfach Reports}, pp.~1347--1406, DOI: \url{https://doi.org/10.4171/owr/2021/26}. The record describes Oberwolfach Reports as documenting mathematical workshops and their open-problem discussions. No individual proposer is identified; the supplied proposed year is 2021.
\item Supplied literature source, labelled ``Oberwolfach report'': \url{https://doi.org/10.4171/owr/2021/26}. Literature checked at: \texttt{2026-08-22}. This exact date belongs to supplied dated curation, not an independently verified literature assessment. That curation labels the inquiry open and describes the comparative questions about classics and their use and reading in Arabic traditions as active historiographical research questions rather than solved propositions. The background calls this dated literature triage derived from supplied OWR curation and requests independent verification of time-sensitive status and cited claims.
\end{itemize}
\end{document}
